\section{Thực nghiệm và đánh giá}
\label{sec:benchmark}
\subsection{Câu hỏi thực nghiệm và phạm vi}
Thực nghiệm xét ba câu hỏi. Thứ nhất, hai biểu diễn TID-set có xuất đúng
toàn bộ itemset và support so với SPMF hay không? Thứ hai, thay Set bằng
BitVector có giảm runtime và Peak RSS trên các loại dữ liệu khác nhau
hay không? Thứ ba, số giao dịch và độ dài giỏ tác động thế nào đến chi phí
khai thác? Correctness và hiệu năng được kiểm tra riêng: output đúng
không tự bảo đảm phương pháp đo công bằng.

Thiết kế gồm bốn nhóm: runtime và output size theo minsup; Peak RSS ở
minsup giữa dải; scalability theo số giao dịch; và dữ liệu tổng hợp tăng
độ dài giỏ. Mỗi nhóm giữ những biến còn lại cố định trong phạm vi thiết kế.
SPMF Eclat~\cite{spmf} là tham chiếu độc lập, không được gọi bên trong
cài đặt Julia của nhóm.

Các số đo hiệu năng của lần chạy được chọn sẽ điền vào các bảng và khung
hình bên dưới. Bản thảo chưa kết luận backend nào thắng trên dataset nào.
Các lập luận về dense/sparse là cơ chế cần kiểm tra, không phải quan sát
đã có của lần đo đang chờ.

\subsection{Dữ liệu, chuẩn hóa và ngưỡng}
Bốn tập dữ liệu từ FIMI~\cite{fimi} bao phủ giao dịch dày và thưa
(Bảng~\ref{tab:data}). Chess/Mushroom có giỏ dài; Retail có nhiều item
nhưng giỏ ngắn; T10I4D100K là dữ liệu tổng hợp. Số giao dịch lớn không
đồng nghĩa output lớn: phân bố đồng xuất hiện và ngưỡng support còn
quyết định số tổ hợp vượt ngưỡng.

\begin{reporttable}
\caption{Thống kê đầu vào benchmark đã chuẩn hóa; $\bar L$ là độ dài giỏ trung bình.}
\label{tab:data}
\small
\begin{tabular}{lrrr}\toprule
CSDL & $N$ & $M$ & $\bar L$\\\midrule
Chess & 3.196 & 75 & 37,0000\\
Mushroom & 8.124 & 119 & 23,0000\\
Retail & 88.162 & 16.470 & 10,3058\\
T10I4D100K & 100.000 & 870 & 10,1023\\\bottomrule
\end{tabular}
\end{reporttable}

Giữ nguyên giao dịch trùng và thứ tự dòng. Chess/Mushroom giữ ID gốc;
Retail/T10I4D100K cộng 1 vào mọi ID để không còn ID 0. Phép đổi tên một-một
này bảo toàn support; Julia và SPMF phải dùng cùng file sau đổi tên.
Không so output có mapping khác nhau chỉ bằng số lượng itemsets.
Script chuẩn bị dữ liệu kiểm tra SHA-256 của file gốc trước khi chuẩn hóa.

Với tỷ lệ $s$ và $N$ giao dịch, ngưỡng đếm là
\begin{equation}
c_{\min}=\lceil sN\rceil.
\label{eq:threshold}
\end{equation}
Dải ngưỡng trong Bảng~\ref{tab:grids} có bảy điểm mỗi dataset, từ cao
xuống thấp. Ngưỡng tương đối giữ ý nghĩa tần suất giữa các mức kích thước
khác nhau, còn ngưỡng đếm thay đổi theo $N$. Nếu thay dải ngưỡng sau pilot,
cần lưu cấu hình cuối cùng trước khi chạy các lượt chính thức.

\begin{reporttable}
\caption{Dải minsup, đơn vị phần trăm; điểm thứ tư dùng cho Peak RSS.}
\label{tab:grids}
\footnotesize
\begin{tabular}{lp{0.65\columnwidth}}\toprule
CSDL & Bảy ngưỡng từ cao xuống thấp\\\midrule
Chess & 95; 92,5; 90; 87,5; 85; 82,5; 80\\
Mushroom & 95; 90; 80; 70; 60; 50; 40\\
Retail & 10; 8; 6; 4; 2; 1; 0,5\\
T10I4D100K & 3; 2,5; 2; 1,5; 1; 0,75; 0,5\\\bottomrule
\end{tabular}
\end{reporttable}

\subsection{Môi trường và phương pháp đo}
Mọi backend chạy trên cùng máy, cùng đầu vào và một luồng Julia. Khi chốt
số liệu, Bảng~\ref{tab:environment} phải được điền từ metadata của chính
lần đo đó. Không lấy cấu hình máy Windows của kết quả cũ điền cho Colab
hoặc máy Linux khác. Ghi riêng phiên bản Julia, Java và hash JAR SPMF;
tên thuật toán trong log không xác định phiên bản thư viện.

\begin{reporttable}
\caption{Môi trường của lần đo được chọn; \pending{} là chờ cập nhật.}
\label{tab:environment}
\small
\begin{tabular}{lp{0.43\columnwidth}}\toprule
Thuộc tính & Giá trị\\\midrule
Hệ điều hành, kiến trúc & \pending\\
CPU, số CPU logic, RAM & \pending\\
Julia, Java, SHA-256 JAR & \pending\\
Ngày đo, commit nguồn & \pending\\
Luồng Julia & 1\\
Số lượt ghi nhận/cấu hình & 5\\\bottomrule
\end{tabular}
\end{reporttable}

\textit{Runtime toàn tiến trình.} Mỗi mẫu tạo tiến trình mới, gồm khởi
động runtime, JIT, đọc dữ liệu, dựng TID-set, mining và ghi output.
Có một lượt chuẩn bị không ghi nhận cho mỗi chương trình; năm lượt đo
sau vẫn là tiến trình mới và vẫn tính JIT. Đây không phải phép đo miner
đã warm-up trong cùng tiến trình. Chuẩn hóa input và so sánh correctness
thực hiện ngoài đồng hồ.

\textit{Thời gian riêng miner.} Scalability và độ dài giao dịch dùng một
lượt warm-up trên đúng CSDL/cùng minsup trong worker, gọi GC trước đo,
rồi đo một lời gọi mine. GC phát sinh bên trong miner vẫn tính vào thời
gian. Phạm vi đo này khảo sát chi phí khai thác; không dùng để so với
runtime SPMF toàn tiến trình ở thí nghiệm khác.

\textit{Bộ nhớ.} Peak RSS là mức bộ nhớ thường trú cao nhất của worker
đến hết mining, đo trước khi sắp xếp/ghi output. Worker warm-up trên toy,
GC, đọc dữ liệu, GC rồi mine. Bộ đếm mức đỉnh không reset, nên đỉnh vẫn
gồm runtime và warm-up. Trên Windows, đại lượng tương ứng lấy từ peak
working set~\cite{windowsrss}. Allocation tích lũy hay kích thước object
Julia không thay thế cho Peak RSS.

Chạy tuần tự và luân phiên thứ tự giữa các lượt để giảm ảnh hưởng ưu tiên
cố định. Báo cáo trung vị và min--max của năm mẫu, không gọi min--max là
khoảng tin cậy. Nếu máy có workload khác, phải ghi điều kiện đó; không
gộp mẫu của hai môi trường thành một median.

\subsection{Correctness và kiểm thử tự động}
Gọi $R$ là output SPMF, $A$ là output Julia, đều được chuẩn hóa thành
ánh xạ tuple item tăng dần sang số đếm. Một cấu hình đạt khi
\begin{equation}
\operatorname{keys}(A)=\operatorname{keys}(R),\quad
\forall X\in R:\ A[X]=R[X].
\end{equation}
Điều kiện này mạnh hơn bằng số lượng: hai output có thể cùng kích thước
nhưng thiếu/thừa itemset hoặc sai support. Tỷ lệ khớp dùng mẫu số là số
itemsets tham chiếu. Nếu tham chiếu rỗng, chỉ xác nhận đạt khi output cũng
rỗng thay vì chia cho 0.

Bảng~\ref{tab:correct} là correctness đã lưu trên năm CSDL đầy đủ. Đây là
đối chiếu với đáp án SPMF có sẵn, không phải số đo hiệu năng đang chờ.
Test không dùng miner của nhóm tạo đáp án cho chính nó.

\begin{reporttable}
\caption{Correctness đã kiểm chứng: cả hai backend khớp toàn bộ itemset/support.}
\label{tab:correct}
\footnotesize
\begin{tabular}{lrrr}\toprule
CSDL & Minsup (\%) & Itemsets SPMF & Khớp\\\midrule
Groceries & 1 & 333 & 100\%\\
Chess & 87,5 & 1.386 & 100\%\\
Mushroom & 70 & 31 & 100\%\\
Retail & 4 & 18 & 100\%\\
T10I4D100K & 1,5 & 237 & 100\%\\\bottomrule
\end{tabular}
\end{reporttable}

Ngoài reference SPMF, test dùng phép vét cạn tập con trên dữ liệu nhỏ và
kiểm tra input không bị sửa tại chỗ. Những ca rỗng, item trùng, giao dịch
trùng, ngưỡng biên và ID sai giúp tách lỗi parser khỏi lỗi DFS. Log package
569/569 assertions trong README bao gồm cả ứng dụng; đây không phải số
dataset hoặc số lượt SPMF. Trước khi cập nhật hiệu năng, cần chạy lại
correctness trên đúng commit và mapping đầu vào.

\subsection{Runtime và output size theo minsup}
Với mỗi dataset và mỗi ngưỡng ở Bảng~\ref{tab:grids}, chạy Set, BitVector
và SPMF năm lần: tổng số mẫu dự kiến là $4\times7\times3\times5=420$.
Mỗi mẫu phải có runtime, số itemsets và trạng thái đối chiếu đầy đủ.
Bảng~\ref{tab:runtime-pending} là khung tổng hợp tại ngưỡng thấp nhất;
CSV vẫn phải lưu cả 28 cấu hình.

\begin{reporttable}
\caption{Runtime toàn tiến trình tại ngưỡng thấp nhất (s); các ô chờ kết quả.}
\label{tab:runtime-pending}
\footnotesize
\begin{tabular}{lrrrrr}\toprule
CSDL & $s$ (\%) & Set & BV & SPMF & Itemsets\\\midrule
Chess & 80 & \pending & \pending & \pending & \pending\\
Mushroom & 40 & \pending & \pending & \pending & \pending\\
Retail & 0,5 & \pending & \pending & \pending & \pending\\
T10I4D100K & 0,5 & \pending & \pending & \pending & \pending\\\bottomrule
\end{tabular}
\end{reporttable}

\begin{reportfigure}
\pendingplot{37mm}{Runtime theo minsup cho bốn CSDL\\Ba đường: Set, BitVector, SPMF\\Trung vị năm lượt và vùng min--max}
\caption{Vị trí đồ thị runtime toàn tiến trình. Minsup đi từ cao xuống thấp; runtime dùng giây. Chưa có đường dữ liệu trong bản thảo.}
\label{fig:runtime-pending}
\end{reportfigure}

Hình~\ref{fig:runtime-pending} cần cho thấy ba đường trên cùng hệ trục cho
mỗi dataset. Nếu dùng thang log, phải ghi rõ. Chi phí khởi động/JIT có thể
làm runtime không tăng đều ở từng điểm dù output tăng; đường toàn tiến
trình không tách được tốc độ riêng phần giao TID-set của hai ngôn ngữ.

\begin{reportfigure}
\pendingplot{33mm}{Số itemsets theo minsup\\Bốn CSDL; đối chiếu ba chương trình\\Đánh dấu các run không hoàn tất}
\caption{Vị trí đồ thị output size; chỉ điền số lượng của run hoàn tất và đã đối chiếu support.}
\label{fig:counts-pending}
\end{reportfigure}

Với input cố định, hạ minsup không thể làm mất itemset đã phổ biến. Đường
output size trong Hình~\ref{fig:counts-pending} phải không giảm khi đi từ
ngưỡng cao xuống thấp. Nếu vi phạm, cần kiểm tra input, mapping hoặc run
chưa hoàn tất trước khi giải thích hiệu năng. Phân tích dense/sparse cần
xét phân bố support và độ dài itemset, không chỉ số giao dịch.

\subsection{Peak RSS tại minsup trung bình}
Chọn điểm thứ tư của mỗi dải: Chess 87,5\%, Mushroom 70\%, Retail 4\%,
T10I4D100K 1,5\%. Mỗi backend đo trong năm worker riêng, tổng cộng
$4\times2\times5=40$ mẫu. Output sau đo phải khớp reference và backend
còn lại. Mẫu lỗi giữ trạng thái/log, không điền 0 vào ô RAM để tính median.

\begin{reporttable}
\caption{Peak RSS trung vị (MiB); $\Delta_M$ chờ tính từ hai median.}
\label{tab:ram-pending}
\small
\begin{tabular}{lrrr}\toprule
CSDL & Set & BV & $\Delta_M$ (\%)\\\midrule
Chess & \pending & \pending & \pending\\
Mushroom & \pending & \pending & \pending\\
Retail & \pending & \pending & \pending\\
T10I4D100K & \pending & \pending & \pending\\\bottomrule
\end{tabular}
\end{reporttable}

\begin{reportfigure}
\pendingplot{32mm}{Peak RSS tại ngưỡng giữa dải\\Hai cột Set/BitVector cho mỗi CSDL\\Median và min--max, đơn vị MiB}
\caption{Vị trí đồ thị bộ nhớ; không thay Peak RSS bằng allocation hoặc RAM Java tự báo.}
\label{fig:ram-pending}
\end{reportfigure}

Với $\widetilde M$ là median của năm đỉnh tiến trình, mức giảm RAM là
\begin{equation}
\Delta_M=100\left(1-\frac{\widetilde M_{BV}}{\widetilde M_{Set}}\right).
\end{equation}
Giá trị dương nghĩa là median BitVector thấp hơn. Công thức không biến
phần chênh thành RAM riêng thuật toán. Không lấy hiệu Peak RSS trước và
sau mining để suy bộ nhớ miner, vì cả hai đều là high-water mark.

Khi phân tích Bảng~\ref{tab:ram-pending} và Hình~\ref{fig:ram-pending},
cần xét số singleton còn lại. Cấp $N$ bit cho item hiếm có thể lãng phí;
lọc sớm giảm số vector nhưng thêm một lượt quét. Muốn tách riêng tác động
của bitset và lọc singleton cần ablation cùng input và cùng miner.

\subsection{Scalability theo số giao dịch}
Retail có 88.162 giao dịch và nhiều item. Dùng
\texttt{MersenneTwister(20261116)} tạo một hoán vị TID, lấy tiền tố
10\%, 25\%, 50\%, 75\%, 100\%, rồi sắp chỉ số theo thứ tự dòng gốc.
Các tập con lồng nhau; mức 100\% khớp file chuẩn hóa ban đầu.
Không nhân bản giao dịch để làm dataset lớn hơn.

Giữ minsup 4\%, tính lại $c_{\min}$ cho từng tập con. Hai backend đo riêng
mine sau warm-up; mỗi mức lặp năm lần, tổng 50 mẫu.
Bảng~\ref{tab:scale-pending} ghi cả kích thước và ngưỡng đếm để không
hiểu nhầm rằng mọi mức dùng cùng min-count.

\begin{reporttable}
\caption{Scalability Retail: thời gian mine sau warm-up (ms), chờ điền.}
\label{tab:scale-pending}
\footnotesize
\begin{tabular}{rrrrr}\toprule
Tỷ lệ (\%) & $N$ & $c_{\min}$ & Set & BV\\\midrule
10 & 8.816 & 353 & \pending & \pending\\
25 & 22.040 & 882 & \pending & \pending\\
50 & 44.081 & 1.764 & \pending & \pending\\
75 & 66.121 & 2.645 & \pending & \pending\\
100 & 88.162 & 3.527 & \pending & \pending\\\bottomrule
\end{tabular}
\end{reporttable}

\begin{reportfigure}
\pendingplot{33mm}{Scalability Retail\\Trục ngang: số giao dịch\\Trục dọc: mine sau warm-up (ms)}
\caption{Vị trí đồ thị scalability với năm mức dữ liệu lồng nhau.}
\label{fig:scale-pending}
\end{reportfigure}

Sau khi có Hình~\ref{fig:scale-pending}, xét tỷ số tăng thời gian so với
tỷ số tăng $N$ và kiểm tra output size ở mỗi mức. Nếu số itemsets đổi,
thời gian tăng có thể do TID-set lớn hơn lẫn nhiều nhánh DFS hơn. Khớp
đường tuyến tính chỉ mô tả miền đã đo, không chứng minh Eclat tuyến tính
với mọi input và minsup.

\subsection{Ảnh hưởng của độ dài giao dịch}
Sinh $N=5.000$ giao dịch trên $M=64$ item, dùng seed 20261116. Mỗi TID có
một hoán vị riêng của $1{:}64$; lấy 4, 8, 12, 16, 20 item đầu để tạo các
mức độ dài. Vì dùng cùng hoán vị cho một TID, giỏ dài chứa giỏ ngắn tương
ứng. Giữ minsup 5\%, nên $c_{\min}=250$ ở mọi mức.

Không loại giao dịch trùng. Các giỏ trong một CSDL có cùng độ dài $L$,
nên độ dài trung bình bằng $L$ và mật độ bằng $L/M$. Thiết kế tách tác
động của độ dài khỏi $N,M$, nhưng support của nhiều tổ hợp vẫn tăng theo
$L$; đó chính là cơ chế cần khảo sát.

\begin{reporttable}
\caption{Dữ liệu tổng hợp và thời gian mine sau warm-up (ms), chờ điền.}
\label{tab:length-pending}
\footnotesize
\begin{tabular}{rrrrr}\toprule
$L$ & Mật độ (\%) & Itemsets & Set & BV\\\midrule
4 & 6,25 & \pending & \pending & \pending\\
8 & 12,50 & \pending & \pending & \pending\\
12 & 18,75 & \pending & \pending & \pending\\
16 & 25,00 & \pending & \pending & \pending\\
20 & 31,25 & \pending & \pending & \pending\\\bottomrule
\end{tabular}
\end{reporttable}

\begin{reportfigure}
\pendingplot{33mm}{Độ dài giao dịch và chi phí khai thác\\Năm mức $L=4,8,12,16,20$\\Runtime và output size trên hai hệ trục riêng}
\caption{Vị trí đồ thị tác động độ dài; giữ $N$, $M$ và minsup cố định.}
\label{fig:length-pending}
\end{reportfigure}

Với chọn đều $L$ item từ $M$, xác suất một cặp cụ thể có mặt trong giỏ là
\begin{equation}
p_2=\frac{L(L-1)}{M(M-1)}.
\end{equation}
Support kỳ vọng $Np_2$ khoảng 163,7 tại $L=12$ và 297,6 tại $L=16$, nằm
hai phía ngưỡng 250. Đây là tính toán từ thiết kế, chưa phải output quan
sát. Nó giải thích vì sao có thể có bước tăng mạnh số itemsets khi giỏ
dài hơn. Khi điền Hình~\ref{fig:length-pending}, cần kiểm tra dữ liệu có
phù hợp cơ chế này và thời gian bị chi phối bởi giao TID-set hay lưu output.

\subsection{Cách tổng hợp và giới hạn kết luận}
Mỗi cấu hình giữ đủ năm mẫu và trạng thái, sau đó mới tính median. Tỷ số
$t_{Set}/t_{BV}>1$ biểu thị BitVector nhanh hơn trong cùng phạm vi đo; các
thời gian phải dương và cùng đơn vị. Nếu run timeout hoặc OOM, không lấy
thời gian giới hạn như thời gian hoàn tất và không đưa output một phần
vào đối chiếu đầy đủ. Số run thành công trên tổng số dự kiến cần được
trình bày cùng kết quả.

Khi cập nhật, thay đồng thời bảng, hình, môi trường, số lần lặp và nhận
xét. Correctness đúng không chứng minh thiết kế đo công bằng; đo nhanh
nhưng output thiếu cũng không phải cải thiện hợp lệ.

Hai hướng tiếp theo là chọn biểu diễn TID theo mật độ và giảm cấp phát
trong DFS. Hướng thứ nhất dùng cấu trúc thưa cho nhánh ít TID và bitset
cho nhánh dày, tính cả chi phí chuyển biểu diễn. Hướng thứ hai có thể
đếm giao trước hoặc dùng lại bộ đệm theo độ sâu; không được ghi đè TID-set
còn dùng ở nhánh khác. Cả hai là đề xuất chưa được coi là đã triển khai.

Một máy, năm lượt và một seed cho mỗi thiết kế chưa đại diện mọi môi
trường. Kết luận chỉ áp dụng cho dải minsup, dataset và phương pháp đã
ghi; không suy từ đường gần tuyến tính rằng thuật toán có cận đa thức
khi phải xuất mọi tập phổ biến.
